% !TEX root = ProjetoPDS.tex
\section{Conclusões}

A comparação entre as quatro aproximações deixou claro que o tipo de filtro
pesa tanto quanto a ordem ou a frequência de corte na qualidade do
resultado. Chebyshev~I e Elíptico convergiram, na busca em grade, para
ordem~1 com $f_c = \SI{600}{\hertz}$ -- um decaimento espectral lento demais
para separar o ruído sem comprometer o restante do sinal. Já o Butterworth
de ordem~3 melhorou esse cenário, mas, com $f_c = \SI{1100}{\hertz}$, acabou
atenuando parte das harmônicas de ordem $6k\pm1$ geradas pelas cargas não
lineares, o que se refletiu no RMSE.

O Chebyshev~II de ordem~3 e $f_c = \SI{2700}{\hertz}$ se destacou justamente
por combinar banda de passagem sem ripple, preservando a fundamental e as
harmônicas mais baixas, com zeros de transmissão na banda de rejeição, que
aceleram a queda do ganho logo após o corte. Essa combinação permitiu
remover o ruído de medição sem distorcer a informação harmônica, que é o que
de fato importa em uma análise de qualidade de energia.

O estudo do custo computacional em função da ordem reforça essa escolha: o
RMSE cai de forma acentuada até a ordem~3 e satura a partir daí, enquanto o
tempo de execução segue crescendo. Não há, portanto, ganho de precisão que
justifique usar uma ordem maior, apenas custo extra em uma aplicação que
poderia, em princípio, rodar em tempo real.

De forma geral, o trabalho mostra que projetar um filtro IIR não se resume a
escolher a aproximação com a transição mais rápida, mas sim a equilibrar
seletividade, fidelidade na banda de passagem e custo de implementação para
o sinal específico em questão.
